\section{Finite-order collision cumulants and spectral damping}
\label{sec:finite-order-damping}
The short-lag argument of Section~\ref{sec:direct-orbit-bounds} does
not estimate a prescribed later coefficient. We now return to the
collision spectral splitting and prove stronger finite-order bounds
before stopping. The point is that its derivatives at zero can be
bounded independently of the smoothing scale, although the analytic
neighborhood still shrinks with that scale. A Taylor polynomial of
fixed, sufficiently high degree then gives real-frequency damping on
a larger band. No analyticity of the unsmoothed twist is asserted.

All expectations in this section use the stationary collision
probability $\nu$. Put $\|u\|_{\mathcal V}=\|u\|_\infty+
\|u\|_{\mathrm{BV}}$. Every order below is fixed before the smoothing
scale, collision length, and radius are varied. Constants may depend
on that order. We require no factorial bound uniform in the order.

\subsection{Connected correlations at any fixed order}
\begin{lemma}[Finite-product decoupling at arbitrary fixed order]
\label{lem:all-finite-product}
For each integer $q\ge2$ there are $C_q,\gamma_q>0$, uniform in $R$,
such that, for $t_1\le\cdots\le t_q$ and $1\le p<q$,
\begin{align}
 &\left|\int\prod_{j=1}^q u_j\circ T_R^{t_j}\dd\nu
 -\left(\int\prod_{j\le p}u_j\circ T_R^{t_j}\dd\nu\right)
  \left(\int\prod_{j>p}u_j\circ T_R^{t_j}\dd\nu\right)\right|
 \notag\\
 &\hspace{12mm}\le C_qe^{-\gamma_q(t_{p+1}-t_p)}
                         \prod_{j=1}^q\|u_j\|_{\mathcal V}.
 \label{eq:all-finite-product}
\end{align}
Negative and repeated times are allowed. The constants can be chosen
simultaneously for orders at most any specified integer $Q$.
\end{lemma}
\begin{proof}
The proof of Lemma~\ref{lem:bv-finite-product} works at each fixed
order, and we give its scale explicitly. By multilinearity normalize
all nonzero norms to one. Smooth every factor at scale $\epsilon$.
Telescoping the products and using invariance bounds the total error
by $C_q\epsilon$. Translate $t_1$ to zero, write the expectation as
the chronological word of smooth multipliers and collision powers,
and replace the power across the selected gap by $\Pi_R=\nu\ell$.
The new word is exactly the product of the two block expectations.
The error is at most $C_q\epsilon^{-2q}\vartheta^{t_{p+1}-t_p}$:
there are $q$ smooth multipliers, all other powers are uniformly
bounded, and the selected complementary power decays exponentially.
For a zero gap use $I-\Pi_R$ with its uniform norm. Taking
$\epsilon=\vartheta^{(t_{p+1}-t_p)/(2q+1)}/4$ proves the assertion.
A finite minimum of the resulting exponents treats orders at most $Q$.
No variation of a long pullback is estimated.
\end{proof}

For bounded real random variables, define the joint cumulant by
\begin{equation}\label{eq:finite-order-cumulant-definition}
 \operatorname{cum}(X_1,\ldots,X_q)
 =\sum_{\pi\in\mathcal P_q}(-1)^{|\pi|-1}(|\pi|-1)!
                    \prod_{B\in\pi}\mathbb E\prod_{j\in B}X_j,
\end{equation}
where $\mathcal P_q$ is the set of partitions of $\{1,\ldots,q\}$.
This is the mixed derivative at zero of the logarithm of the joint
moment generating function. In particular it is multilinear and
vanishes when the variables split into two nonempty independent
families. The latter assertion follows because the joint generating
function factors and its logarithm is a sum of two functions in
separate sets of variables. This elementary finite-order use of
cumulants is distinct from an all-orders normal-approximation theorem;
see also \cite{DJS} for the general framework.

\begin{proposition}[Absolute summability of fixed-order connected correlations]
\label{prop:connected-collision-summability}
Let $u$ be real and $\|u\|_{\mathcal V}\le B$. For each $q\ge2$,
\begin{equation}\label{eq:connected-collision-summability}
 \sum_{\mathbf k\in\Z^{q-1}}
 (1+\operatorname{diam}\{0,k_2,\ldots,k_q\})
 \left|\operatorname{cum}(u,u\circ T_R^{k_2},\ldots,
                                      u\circ T_R^{k_q})\right|
 \le C_{q,B}.
\end{equation}
Consequently the absolutely convergent sum without the diameter
factor, denoted by $\mathfrak c_q(R,u)$, satisfies
\begin{equation}\label{eq:cumulant-linear-length}
 \left|\operatorname{cum}_q(S_{m,R}u)-m\mathfrak c_q(R,u)\right|
                \le C_{q,B}\qquad(m\ge1),
 \qquad |\mathfrak c_q(R,u)|\le C_{q,B}.
\end{equation}
Both statements are uniform in $R$ and in the displayed class of $u$.
\end{proposition}
\begin{proof}
Order any $q$ times and split them at a largest consecutive gap $d$.
Compare their joint law with a law in which the entire left block and
the entire right block are independent, but the two block marginals
are unchanged. For every subset meeting both blocks, its moment
changes by at most $C_{q,B}e^{-\gamma_qd}$ by
Lemma~\ref{lem:all-finite-product}; the gap in that subset is at least
$d$. All other subset moments are unchanged. The finite partition
polynomial \eqref{eq:finite-order-cumulant-definition}, whose factors
are bounded, therefore changes by at most the same order. Its value
under the independent-block law is zero. Thus the original joint
cumulant is bounded by $C_{q,B}e^{-\gamma_qd}$. If all times coincide,
the bounded partition formula gives this estimate with $d=0$.

For anchored times $0,k_2,\ldots,k_q$, their diameter $D$ obeys
$d\ge D/(q-1)$. There are at most $(2D+1)^{q-1}$ tuples with all
coordinates of absolute value at most $D$. Exponential decay therefore
sums with the additional factor $1+D$, proving
\eqref{eq:connected-collision-summability}.

By multilinearity and stationarity, the exact finite sum is
\[
 \operatorname{cum}_q(S_{m,R}u)
 =\sum_{\mathbf k\in\Z^{q-1}}(m-D(\mathbf k))_+
   \operatorname{cum}(u,u\circ T_R^{k_2},\ldots,u\circ T_R^{k_q}).
\]
Indeed, after fixing the first index, the remaining indices have these
relative offsets; exactly $(m-D)_+$ translates fit in the interval of
$m$ collisions. The difference between this coefficient and $m$ has
absolute value at most $D$. Applying the first assertion proves
\eqref{eq:cumulant-linear-length}. This also treats repeated indices
without an additional combinatorial multiplicity.
\end{proof}

\subsection{Uniform derivatives of the smoothed eigenvalue}
For a real unit vector $\xi\in\R^4$ let
\[
 u_{R,\delta,\xi}=\xi\cdot h_{R,\delta},\qquad
 \psi_{R,\delta,\xi}(t)=\log\lambda_{R,\delta}(t\xi),
\]
using the branch of the logarithm fixed in
Lemma~\ref{lem:smooth-collision-expansion}. The supremum and BV norms
of $u_{R,\delta,\xi}$ are bounded uniformly, and its mean is zero.

\begin{theorem}[A smoothing-independent finite spectral jet]
\label{thm:finite-spectral-jet}
For each fixed $Q\ge3$ there is $C_Q$ such that
\begin{align}
 \psi'(0)&=0,\qquad
 \psi''(0)=-\xi^{\mathsf T}\Gamma_{R,\delta}\xi,\notag\\
 \psi^{(q)}(0)&=i^q\mathfrak c_q(R,u_{R,\delta,\xi}),\qquad
 |\psi^{(q)}(0)|\le C_Q\quad(2\le q\le Q).
 \label{eq:finite-spectral-jet}
\end{align}
In particular, on $|t|\le a\delta^2$, after decreasing $a>0$,
\begin{equation}\label{eq:finite-spectral-Taylor}
 \psi(t)=-\tfrac12t^2\xi^{\mathsf T}\Gamma_{R,\delta}\xi
           +\sum_{q=3}^{Q}\frac{\psi^{(q)}(0)}{q!}t^q
           +O_Q(\delta^{-2(Q+1)}|t|^{Q+1}).
\end{equation}
Every constant is uniform in $R,\delta,\xi$. The Taylor remainder still
uses the shrinking complex analytic radius; it is not scale independent.
\end{theorem}
\begin{proof}
Fix $R,\delta,\xi$ temporarily. The stationary spectral pairing is
\[
 F_m(t):=\int e^{itS_{m,R}u_{R,\delta,\xi}}\dd\nu
               =e^{m\psi(t)}A(t)+E_m(t),\qquad A(t)=\ell(\Pi(t\xi)\nu).
\]
Here $A(0)=1$, and the complementary term has norm at most $C\rho^m$
on the given complex neighborhood. On a smaller neighborhood,
$|A(t)|\ge1/2$ and $|e^{\psi(t)}|\ge\rho_1>\rho$. Hence the relative
error is analytic and bounded by $C(\rho/\rho_1)^m$. For large $m$,
choose its logarithm near one. Cauchy's inequalities on a still smaller
neighborhood give, for each fixed $q$,
\[
 \frac1m(\log F_m)^{(q)}(0)\longrightarrow\psi^{(q)}(0).
\]
The derivatives of $\log A$ divided by $m$ vanish. All these
identification steps are made at fixed $\delta$; their constants need
not be uniform as $\delta\downarrow0$.

Boundedness of the finite collision sum and the cumulant definition
give $(\log F_m)^{(q)}(0)=i^q\operatorname{cum}_q(S_m u)$.
Proposition~\ref{prop:connected-collision-summability} identifies the
limit and bounds it independently of $\delta,R,\xi$.
The first derivative vanishes by centering; the second is the collision
Green--Kubo variance, also identified in
Lemma~\ref{lem:smooth-collision-expansion}.

Finally the analytic logarithm is uniformly bounded on a complex disk
of radius a fixed multiple of $\delta^2$, as in the proof of that
lemma. Cauchy's estimate for the Taylor tail on a concentric smaller
disk is $C_Q\delta^{-2(Q+1)}|t|^{Q+1}$. This proves
\eqref{eq:finite-spectral-Taylor}. Only the finitely many coefficients,
not the analytic disk or an infinite cumulant series, have gained
smoothing-independent bounds.
\end{proof}

\begin{corollary}[Real-frequency collision damping]
\label{cor:finite-jet-damping}
For each fixed $Q\ge3$ there are $c,C,a,\delta_0>0$ such that
$0<\delta\le\delta_0$ and
\begin{equation}\label{eq:finite-jet-damping-domain}
 |z|\le a\delta^2,\qquad
 \delta^{-2(Q+1)}|z|^{Q-1}\le a
\end{equation}
imply, on every local collision space,
\begin{equation}\label{eq:finite-jet-damping}
 |\lambda_{R,\delta}(z)|\le e^{-c|z|^2},\qquad
 \|\mathcal L_{R,\delta,z}^{m}\|\le Ce^{-cm|z|^2}
                \quad(m\ge0,\ z\in\R^4).
\end{equation}
\end{corollary}
\begin{proof}
Theorem~\ref{thm:joint-uniform-nondegeneracy} and
\eqref{eq:collision-covariance-limit} give
$\Gamma_{R,\delta}\ge c_1I_4$ for all sufficiently small $\delta$,
uniformly in $R$. Write $z=t\xi$. The finite sum of degree at least
three in \eqref{eq:finite-spectral-Taylor} is bounded by $C_Q|t|^3$
for $|t|\le1$. Its remainder is at most $C_Qa|t|^2$ under
\eqref{eq:finite-jet-damping-domain}. Decreasing $a$ and $\delta_0$
therefore gives $\operatorname{Re}\psi(t)\le-c_2|t|^2$.
The bounded spectral projection and the complementary power estimate
then give
$\|\mathcal L_z^m\|\le C(e^{-c_2m|z|^2}+\rho^m)$.
Decrease the constant in the exponent so that $\rho^m$ is dominated
by the same bound on the stated small real-frequency region. This
also holds for $m=0$ after increasing $C$.
\end{proof}

The estimate is a power bound for the \emph{smoothed collision}
operator on its previously constructed local distribution space.
It is not an operator-norm decay assertion for the unitary twisted
return Koopman operator. The next section removes smoothing and
stops at the actual prescribed return, keeping those errors explicit.
